\documentclass[aspectratio=169,11pt]{beamer}

\usetheme{Madrid}
\usecolortheme{seahorse}
\setbeamertemplate{navigation symbols}{}
\setbeamertemplate{footline}[frame number]

\usepackage{amsmath,amssymb}
\usepackage{booktabs}
\usepackage{multirow}
\usepackage{graphicx}
\usepackage{natbib}
\usepackage{tikz}

% Clean fonts
\usefonttheme{professionalfonts}

\title[Information \& Parental Investment]{Why Does the Channel Matter?\\[4pt]
\small Separating Attention and Belief Mechanisms\\in Parental Information Interventions}
\author{Jordi Torres Vallverd\'{u}}
\institute{MRes TSE --- Behavioral Development Economics}
\date{March 2026}

\begin{document}

% ============================================================
% TITLE
% ============================================================
\begin{frame}
\titlepage
\end{frame}

% ============================================================
% OUTLINE
% ============================================================
\begin{frame}{Outline}
\tableofcontents
\end{frame}

% ============================================================
\section{Introduction}
% ============================================================

\begin{frame}{Information Problems in Education}

\textbf{Students and parents make decisions under imperfect information}

\begin{itemize}
    \item Students hold biased beliefs about returns to education and own ability
    \begin{itemize}
        \item Major choice and labor market outcomes \citep{arcidiacono2024recovering,wiswall2015determinants}
        \item Persistence and dropout decisions \citep{stinebrickner2012learning,arcidiacono2016college}
        \item College application and matching \citep{larroucau2024application}
    \end{itemize}

    \medskip

    \item \textbf{Parents are also misinformed} about their children's performance
    \begin{itemize}
        \item Mean belief--performance gap $>$ 1 SD in some contexts \citep{dizonross2019parents}
        \item Biases concentrated among low-SES households
    \end{itemize}

    \medskip

    \item Parental investments respond to information $\Rightarrow$ misperceptions distort human capital accumulation
\end{itemize}

\end{frame}

% ============================================================
\section{Experimental Evidence}
% ============================================================

\begin{frame}{Information Treatments to Parents}

\textbf{Providing information to parents improves outcomes}

\begin{itemize}
    \item \citet{dizonross2019parents}: Correcting beliefs about child ability in Malawi
    \begin{itemize}
        \item Parents reallocate investments toward lower-performing children
        \item Beliefs are systematically upward-biased
    \end{itemize}

    \bigskip

    \item \citet{bergman2021information}: SMS-based grade notifications in the US
    \begin{itemize}
        \item Parents receive same info via SMS vs.\ standard report cards
        \item Result: 0.19 SD GPA increase at very low cost
        \item $\sim$59\% from belief revision, $\sim$41\% from reduced monitoring costs
    \end{itemize}

    \bigskip

    \item \textbf{The puzzle:} the \emph{channel} of delivery matters, not just content
\end{itemize}

\end{frame}

% ============================================================
\section{Mechanisms}
% ============================================================

\begin{frame}{Why Does the Channel Matter? Possible Mechanisms}

Standard economics: only content should matter $\Rightarrow$ \textbf{three behavioral mechanisms}

\bigskip

\begin{enumerate}
    \item \textbf{Biased beliefs \& asymmetric updating}
    \begin{itemize}
        \item Parents hold upward-biased priors; update asymmetrically
        \item More information $\Rightarrow$ belief correction
    \end{itemize}

    \medskip

    \item \textbf{Limited attention \& salience} \textcolor{blue}{$\leftarrow$ \textit{main focus}}
    \begin{itemize}
        \item Under cognitive load, ``child performance'' drops out of active decision model \citep{gabaix2014sparsity}
        \item SMS creates high salience contrast vs.\ report card among paperwork \citep{bordalo2012salience}
    \end{itemize}

    \medskip

    \item \textbf{Present bias \& procrastination}
    \begin{itemize}
        \item Parents intend to engage but delay; SMS reduces action delay
    \end{itemize}
\end{enumerate}

\bigskip
\textbf{Problem:} existing designs bundle content and channel $\Rightarrow$ cannot separate mechanisms
\end{frame}

% ============================================================
\section{Research Proposal}
% ============================================================

\begin{frame}{Research Questions}

\textbf{Goal:} Disentangle the mechanisms behind parental information treatments

\bigskip

\begin{enumerate}
    \item Does changing the \emph{channel} alone (holding content fixed) affect parental behavior?\\
    $\Rightarrow$ Attention / salience mechanism

    \bigskip

    \item Does providing richer \emph{content} (holding channel fixed) affect behavior?\\
    $\Rightarrow$ Belief correction mechanism

    \bigskip

    \item Are channel and content \emph{complements} or \emph{substitutes}?\\
    $\Rightarrow$ Interaction between attention and beliefs

    \bigskip

    \item Does \emph{frequency} of contact matter beyond content and channel?\\
    $\Rightarrow$ Sustained attention vs.\ one-shot correction
\end{enumerate}

\end{frame}

% ============================================================
\section{Model}
% ============================================================

\begin{frame}{Model: Production Function \& Belief Updating}

\textbf{Grade production:} Student $i$ in school $e$:
\[
Y_{i} = Y(a_{i},\, e_{i}^{p},\, e^{s}) + \varepsilon_{i}
\]
\begin{itemize}
    \item $a_i$: student ability \quad $e_i^p$: parental effort \quad $e^s$: school effort
\end{itemize}

\bigskip

\textbf{Belief updating:} Parents hold prior $\hat{a}_{i}^{\text{pre}}$, receive signal $\sigma_i$:
\[
\hat{a}_{i}^{\text{post}} = \hat{a}_{i}^{\text{pre}}
  + \underbrace{A(\kappa_i)}_{\text{attention}} \cdot \underbrace{\lambda(\sigma_i, \hat{a}_{i}^{\text{pre}})}_{\text{salience}}
  \cdot (\sigma_i - \hat{a}_{i}^{\text{pre}})
\]

\begin{itemize}
    \item $A(\kappa_i) \in [0,1]$: attention weight \citep{gabaix2014sparsity}
    \begin{itemize}
        \item High cognitive load $\kappa_i$ $\Rightarrow$ low attention; SMS lowers $\kappa_i$
    \end{itemize}
    \item $\lambda(\cdot) \in [0,1]$: salience weight \citep{bordalo2012salience}
    \begin{itemize}
        \item Signals contrasting with priors get more weight
    \end{itemize}
\end{itemize}

\end{frame}

% ============================================================

\begin{frame}{Model: Effort Choice}

\textbf{Parental effort:} Given updated beliefs, parent maximizes:
\[
\max_{e_i^p \geq 0} \; \ln Y\!\left(a_i,\, e_i^p,\, e^s\right)\Big|_{\hat{a}_{i}^{\text{post}}} - C^p(e_i^p,\, z_i)
\]

\begin{itemize}
    \item $C^p(e_i^p, z_i)$: cost of effort, heterogeneous in parent type $z_i$ (SES, time)
    \item Allows credit constraints to enter through cost function
\end{itemize}

\bigskip

\textbf{Key distinction:}
\begin{itemize}
    \item $A(\kappa_i)$: determines whether parent \emph{considers} the information at all
    \item $\lambda(\cdot)$: determines how much weight the signal gets \emph{conditional on attention}
\end{itemize}

\bigskip

$\Rightarrow$ Both needed: a parent can attend to a signal yet underweight it (or vice versa)

\end{frame}

% ============================================================
\section{Experimental Design}
% ============================================================

\begin{frame}{The RCT: Five Treatment Arms}

\textbf{Setting:} Primary/secondary schools with digital gradebooks but low baseline parent--school communication. Randomization stratified by school characteristics.

\bigskip

\begin{table}[h]
\centering
\small
\begin{tabular}{@{}llcc@{}}
\toprule
\textbf{Arm} & \textbf{Description} & \textbf{Channel} & \textbf{Content} \\
\midrule
\textbf{C}  & Standard report cards (status quo) & Paper & Standard \\
\textbf{T1} & Same info, delivered via SMS/WhatsApp & \textcolor{blue}{SMS} & Standard \\
\textbf{T2} & Detailed info (grades, ranking, absences) & Paper & \textcolor{blue}{Detailed} \\
\textbf{T3} & Detailed info via SMS (replicates Bergman) & \textcolor{blue}{SMS} & \textcolor{blue}{Detailed} \\
\textbf{T4} & Same as T3, double frequency (weekly) & \textcolor{blue}{SMS} & \textcolor{blue}{Detailed+} \\
\bottomrule
\end{tabular}
\end{table}

\bigskip

\begin{itemize}
    \item T1 vs.\ C: isolates \emph{channel} effect (attention)
    \item T2 vs.\ C: isolates \emph{content} effect (beliefs)
    \item T3 vs.\ C: joint effect (replication of Bergman)
    \item T4 vs.\ T3: \emph{frequency} / sustained attention
\end{itemize}

\end{frame}

% ============================================================
\section{Identification}
% ============================================================

\begin{frame}{Identification Strategy}

\textbf{$2 \times 2$ factorial design} enables clean decomposition:

\medskip

\begin{table}[h]
\centering
\begin{tabular}{@{}lcc@{}}
\toprule
 & \textbf{Standard content} & \textbf{Detailed content} \\
\midrule
\textbf{Paper channel} & C & T2 \\
\textbf{SMS channel}   & T1 & T3 \\
\bottomrule
\end{tabular}
\end{table}

\medskip

\textbf{\textcolor{blue}{revise. I think this is a bit off}}

\textbf{Estimating equation:}
\[
Y_i = \alpha + \beta_1 \cdot \mathbf{1}[\text{SMS}]_i + \beta_2 \cdot \mathbf{1}[\text{Detailed}]_i + \beta_3 \cdot (\text{SMS} \times \text{Detailed})_i + X_i'\gamma + \epsilon_i
\]

\medskip

\begin{itemize}
    \item $\beta_1$: attention effect $A$ (channel changes, content fixed)
    \item $\beta_2$: belief-correction effect through $\lambda \cdot (\sigma - \hat{a})$
    \item $\beta_3$: complementarity between attention and beliefs
\end{itemize}

\medskip

\textbf{Additional:} Elicit baseline beliefs $\hat{a}_i^{\text{pre}}$ to separate biased priors from distorted updating. T4 adds frequency dimension.

\end{frame}

% ============================================================
\section{Power Calculations}
% ============================================================

\begin{frame}{Power Calculations}

\textbf{Benchmark:} \citet{bergman2021information} finds 0.19 SD effect (joint treatment)

\medskip

\textbf{Assumptions:}
\begin{itemize}
    \item MDE = 0.10 SD (individual mechanisms smaller than joint effect)
    \item Significance level $\alpha = 0.05$, power $= 0.80$
    \item ICC $\rho = 0.05$ (school-level clustering), 20 students per cluster
\end{itemize}

\medskip

\textbf{Required sample (per arm):}
\[
n = \frac{(z_{1-\alpha/2} + z_{\text{power}})^2 \cdot 2\sigma^2 \cdot \text{DEFF}}{\text{MDE}^2}
\approx \frac{(1.96 + 0.84)^2 \cdot 2 \cdot 1.95}{0.01} \approx 3{,}060
\]

\smallskip

where $\text{DEFF} = 1 + (20-1) \times 0.05 = 1.95$ (design effect for clustering)
\textbf{\textcolor{blue}{unfeasible}}
\medskip

\begin{itemize}
    \item $\approx$ 3,000 students per arm $\times$ 5 arms $=$ \textbf{15,000 students total}
    \item $\approx$ 750 school clusters (150 per arm)
    \item Stratification by school SES reduces variance
\end{itemize}

\end{frame}

% ============================================================
\section{Conclusion}
% ============================================================

\begin{frame}{Conclusion \& Policy Implications}

\textbf{Summary:}
\begin{itemize}
    \item Existing evidence: information treatments improve educational outcomes
    \item Open question: \emph{why}---attention, beliefs, or their interaction?
    \item Proposed RCT with factorial design separates these mechanisms
\end{itemize}

\bigskip

\textbf{Policy implications depend on the mechanism:}
\begin{itemize}
    \item If $\beta_1 \gg \beta_2$: channel/salience drives the effect $\Rightarrow$ invest in \textbf{nudge design}
    \item If $\beta_2 \gg \beta_1$: content/beliefs drive the effect $\Rightarrow$ \textbf{one-shot information} suffices
    \item If $\beta_3 > 0$: attention and beliefs are complements $\Rightarrow$ \textbf{both needed together}
\end{itemize}

\bigskip

\textbf{Contribution:}
\begin{itemize}
    \item From reduced-form to structural understanding of why information works
    \item Low-cost, scalable interventions with large potential returns
    \item Framework generalizable to other informational settings in education
\end{itemize}

\end{frame}

% ============================================================
% REFERENCES
% ============================================================

\begin{frame}[allowframebreaks]{References}
\footnotesize
\bibliographystyle{plainnat}
\bibliography{references_v2}
\end{frame}

\end{document}
